% 2009 数学一 Q03：原题函数图像的矢量重绘。
% 依据原卷扫描图按坐标关系重建：[-1,0) 上为正的常值段；
% (0,2] 上为单调上升曲线，且在 x=1 处过 x 轴；[2,3] 上 f(x)=0。

% 坐标轴
\draw[->, dashed, line width=0.8pt] (-1.55, 0) -- (3.55, 0) node[below=1pt] {$x$};
\draw[->, dashed, line width=0.8pt] (0, -1.25) -- (0, 2.15) node[left=1pt] {$f(x)$};
\node[below left=1pt] at (0,0) {$O$};

% 刻度
\node[below] at (-1,0) {$-1$};
\node[below] at (1,0) {$1$};
\node[below] at (2,0) {$2$};
\node[below] at (3,0) {$3$};
\node[left] at (0,1) {$1$};

% [-1,0)：原图为水平常值段 f=1
\draw[line width=1.1pt] (-1,1) -- (-0.07,1);
\draw[line width=0.9pt, fill=themebg] (0,1) circle (1.7pt);

% (0,2]：按原卷扫描图形状作平滑矢量重建。
% 取拟合曲线 f(x)=0.55x^2+0.05x-0.60；它严格满足 f(1)=0，
% 并保持原图的“先负后正、单调上升”关系。
\draw[line width=1.1pt, domain=0:2, samples=70]
  plot (\x,{0.55*\x*\x+0.05*\x-0.60});

% [2,3]：原图落在 x 轴上
\draw[line width=1.1pt] (2,0) -- (3,0);
